% Apéndice A - Manual del administrador.
% Contenido pendiente. Insumos previstos: docs/tutorial-profesor.md
% (instalación rápida), docs/environment.md, claude/decisiones.md
% (filosofía resumida), referencia de comandos.
%
% Va dirigido al administrador / superusuario que se enfrenta a
% UnivOrch por primera vez. Debe entender la filosofía, instalarlo
% en su infraestructura, configurar hipervisores reales, gestionar
% usuarios (cuando aplique) y mantenerlo.

\umaappendix{Manual del administrador}
\label{ap:manual-admin}

% TODO: redactar. Estructura sugerida:
% A.1 Filosofía resumida (1-2 páginas)
%   - Qué es UnivOrch, modelo de árbol, herencia, dos ejes de estado.
% A.2 Instalación rápida (basada en tutorial-profesor.md)
%   - Requisitos, instalador de una línea, primera arranque.
% A.3 Instalación en producción
%   - Configuración del docker-compose para producción real.
%   - Volumen persistente y respaldo del fichero TinyDB.
%   - Despliegue tras proxy inverso (sin TLS en el contenedor).
% A.4 Configuración de hipervisores reales
%   - VMware: credenciales, certificados, plantillas base.
%   - Proxmox: API token, nodos.
%   - Mock para pruebas.
% A.5 Carga de definiciones
%   - Sintaxis del YAML: carpetas, descriptores, recursos (define/use).
%   - Idempotencia del load, qué pasa cuando se carga dos veces.
% A.6 Operación habitual
%   - deploy/start/stop/undeploy desde CLI.
%   - Lectura de Jobs y diagnóstico.
% A.7 Mantenimiento
%   - Logs del daemon (dónde y cómo consultar).
%   - Histórico de operaciones desde la BD.
%   - Copias de seguridad del volumen.
%   - Actualización de versión (docker compose pull + restart).

Pendiente de redactar. Se incluirá en una versión posterior.
